\documentclass[conference]{IEEEtran}

\usepackage{cite}
\usepackage{amsmath,amssymb,amsfonts}
\usepackage{graphicx}
\usepackage{textcomp}
\usepackage{xcolor}
\usepackage{booktabs}
\usepackage{tagging}
\usepackage{url}

\usepackage{sc26repro}

% Uncomment/comment the following usetag lines
% if you would like to get an explanation or an example.
% Don't forget to comment them for the final submission.
%\usetag{explanation}
%\usetag{example}

\begin{document}

\twocolumn[%
{\begin{center}
\Huge
Appendix: Artifact Description/Artifact Evaluation
\end{center}}
]

%%%%%%%%%%%%%%%%%%%%%%%%%%%%%%%%%%%%%%%%%%%%%%%%%%%%%
%  AD Appendix
%%%%%%%%%%%%%%%%%%%%%%%%%%%%%%%%%%%%%%%%%%%%%%%%%%%%%

\appendixAD

\section{Overview of Contributions and Artifacts}

\subsection{Paper's Main Contributions}

\begin{description}
\item[$C_1$] Science-aware retrieval operators: deterministic SI dimensional conversion and formula normalization implemented as first-class retrieval branches alongside BM25 and dense vector search, with correctness guaranteed by construction through a 58-unit registry across 12 SI domains.
\item[$C_2$] Corpus-adaptive search orchestration: a profiling layer that inspects each namespace before querying (document count, measurement count, metadata density) and selects search strategies accordingly, activating science-aware operators only where extractable measurements exist.
\item[$C_3$] Federated scientific search: a connector architecture spanning heterogeneous HPC backends (filesystems, object stores, HDF5, NetCDF, IOWarp tiered blob storage, external catalogs) with per-connector capability negotiation, and an agentic multi-hop refinement loop that expands, narrows, or pivots queries based on intermediate results.
\end{description}


\subsection{Computational Artifacts}

\begin{description}
\item[$A_1$] \url{https://github.com/SIslamMun/clio-search}
\item[$A_2$] \url{https://github.com/SIslamMun/clio-search/tree/main/eval/eval_final/code/laptop}
\item[$A_3$] \url{https://github.com/SIslamMun/clio-search/tree/main/eval/eval_final/code/delta}
\end{description}

\begin{center}
\begin{tabular}{rll}
\toprule
Artifact ID  &  Contributions &  Related \\
             &  Supported     &  Paper Elements \\
\midrule
$A_1$   &  $C_1$, $C_2$ & Tables IV, VI--VIII \\
        &               & Figure 3 \\
\midrule
$A_2$   &  $C_1$, $C_2$ & Table V \\
        &               & Figure 2 \\
\midrule
$A_3$   &  $C_2$, $C_3$ & Tables IX, X \\
        &               & Figures 4, 5 \\
\bottomrule
\end{tabular}
\end{center}


%%%%%%%%%%%%%%%%%%%%%%%%%%%%%%%%%%%%%%%%%%%%%%%%%%%%%%%%
\section{Artifact Identification}
%%%%%%%%%%%%%%%%%%%%%%%%%%%%%%%%%%%%%%%%%%%%%%%%%%%%%%%%

%%%%%%%%%%%%%%%%%%%%%%%%%%%%%%%%%%%%%%%%%%%%%%%%%%%%%%%%
%% A1 -- Deterministic retrieval experiments (no API)
%%%%%%%%%%%%%%%%%%%%%%%%%%%%%%%%%%%%%%%%%%%%%%%%%%%%%%%%

\newartifact

\artrel

Artifact $A_1$ contains the CLIO Search implementation (13{,}100 lines, 60 modules, 283 tests), the 210-document controlled benchmark, and the 288-document NOAA GHCN-Daily real-world sample. It supports $C_1$ (science-aware operators) and $C_2$ (corpus-adaptive orchestration) through three deterministic single-node retrieval experiments that require no external API access.

\artexp

Reproduces Table~IV (benchmark corpus composition), Table~VI (cross-unit retrieval at $K{=}5$ with CLIO achieving P@5 of 0.99 versus 0.40 for BM25), Table~VII (NOAA GHCN-Daily validation with CLIO at 0.700~P@5 versus 0.375 for the hybrid baseline), Table~VIII (ablation showing science operators as the dominant gain: P@5 jumps from 0.36 to 0.99 once operators are added), and Figure~3 (cross-unit P@5 bar chart).

\arttime

Artifact Setup: 15 minutes (\texttt{uv sync} plus one-time embedding model download).
Artifact Execution: approximately 15 minutes --- controlled cross-unit benchmark $\sim$5 min, NOAA real-world validation $\sim$5 min, ablation $\sim$5 min.
Artifact Analysis: 5 minutes (plot generation).
Total: approximately 35 minutes.

\artin

\artinpart{Hardware}

Any x86\_64 Linux system with 16~GB RAM and 20~GB free disk. No GPU required. No network required after the initial model download.

\artinpart{Software}

\begin{itemize}
\item Python 3.11 or newer \url{https://www.python.org/}
\item \texttt{uv} package manager \url{https://github.com/astral-sh/uv}
\item DuckDB 1.1.0 (installed via Python dependencies)
\item all-MiniLM-L6-v2 sentence-transformer (downloaded on first run)
\end{itemize}

\artinpart{Datasets / Inputs}

The 210-document controlled benchmark (Table~IV) is generated by \texttt{code/benchmarks/generate\_corpus\_v2.py} and also ships pre-generated under \texttt{code/benchmarks/corpus\_v2/} (seven domain subdirectories: materials science, atmospheric science, fluid dynamics, chemistry, HPC simulation, cross-domain mixed, negative controls). Ground truth is generated programmatically: a document is relevant to a query if and only if it contains a measurement within the query's numeric range after SI conversion. The NOAA GHCN-Daily real-world corpus (1{,}728 station-month files; 288 are sampled for Table~VII) is included in the repository under \texttt{code/benchmarks/corpus\_real/}.

\artinpart{Installation and Deployment}

\begin{verbatim}
git clone https://github.com/SIslamMun/clio-search
cd clio-search/code
uv sync --all-extras --dev
uv run pytest tests/ -v
\end{verbatim}

\artcomp

$T_1$: generate the 210-document synthetic corpus (\texttt{generate\_corpus\_v2.py}) and index it into DuckDB with SI canonicalization;
$T_2$: run \texttt{eval/eval\_final/code/laptop/L4\_numconq\_benchmark.py} to execute 80 queries across five methods (BM25, Dense, Hybrid, String norm., CLIO) and record P@5, R@5, F1@5, MRR (Table~VI);
$T_3$: run \texttt{code/benchmarks/evaluate\_real.py} over the 288-document NOAA sample drawn from \texttt{corpus\_real/} (Table~VII);
$T_4$: run \texttt{eval/eval\_final/code/laptop/L3\_si\_unit\_cross\_prefix.py} to produce the cumulative-components ablation (Table~VIII);
$T_5$: run \texttt{eval/eval\_final/code/generate\_plots.py} to emit Figure~3 from the Table~VI JSON.

Dependencies: $T_1 \rightarrow \{T_2, T_3, T_4\} \rightarrow T_5$; tasks $T_2$, $T_3$, $T_4$ are independent once $T_1$ is complete.

Parameters: BM25 uses $k_1{=}1.2$, $b{=}0.75$; the dense branch uses all-MiniLM-L6-v2 (384-dim) with cosine similarity; $K{=}5$ for all retrieval metrics; ablation follows the cumulative-component staging A (BM25) $\rightarrow$ B (+dense) $\rightarrow$ C (+science operators) $\rightarrow$ D (full pipeline).

\artout

Raw results are written as JSON under \texttt{eval/eval\_final/outputs/}: \texttt{L4\_numconq\_benchmark.json} (Table~VI), the NOAA run (Table~VII), and \texttt{L3\_si\_unit\_cross\_prefix.json} (Table~VIII). Figure~3 (\texttt{plots/fig\_crossunit\_p5\_bar.png}) is generated from the Table~VI JSON via \texttt{generate\_plots.py}.


%%%%%%%%%%%%%%%%%%%%%%%%%%%%%%%%%%%%%%%%%%%%%%%%%%%%%%%%
%% A2 -- Agentic three-path efficiency (Claude API)
%%%%%%%%%%%%%%%%%%%%%%%%%%%%%%%%%%%%%%%%%%%%%%%%%%%%%%%%

\newartifact

\artrel

Artifact $A_2$ contains the agentic three-path evaluation: 10~scientific discovery queries issued through (1) a Claude Native Agent with web search and sub-agents, (2) an agent with NDP-MCP catalog access, and (3) an agent backed by CLIO Search indexing 341 NDP datasets. It supports $C_1$ (science-aware operators) and $C_2$ (corpus-adaptive orchestration) by demonstrating that CLIO's local deterministic operators reduce LLM token consumption and wall time compared to raw catalog access.

\artexp

Reproduces Table~V (CLIO achieves 52.08\% token reduction and 91.27\% wall-time reduction versus the Claude Native Agent, with 10/10 correct answers on the same queries) and Figure~2 (token and wall-time bar chart).

\arttime

Artifact Setup: 5 minutes after $A_1$ is installed (export \texttt{ANTHROPIC\_API\_KEY}).
Artifact Execution: approximately 75 minutes --- Claude Native Agent path $\sim$51~min (spawns 9~sub-agents), NDP-MCP path $\sim$17~min, CLIO path $\sim$4~min; 10 queries each, single repetition.
Artifact Analysis: 5 minutes (plot generation).
Total: approximately 85 minutes.

\artin

\artinpart{Hardware}

Same as $A_1$ (x86\_64 Linux, 16~GB RAM, 20~GB disk), plus outbound HTTPS to the Claude API endpoint.

\artinpart{Software}

\begin{itemize}
\item All of $A_1$'s software stack
\item Claude Agent SDK
\item \texttt{claude-sonnet-4-20250514} API access (paid). Required to reproduce Table~V; a deterministic fallback lets the CLIO path run without it but does not reproduce the agentic token/time comparison.
\end{itemize}

\artinpart{Datasets / Inputs}

Ten scientific discovery queries are embedded in \texttt{L1\_three\_path.py}. The National Data Platform (NDP) catalog (341 datasets) is fetched at runtime and indexed locally by the CLIO path.

\artinpart{Installation and Deployment}

After completing $A_1$'s \texttt{uv sync}:
\begin{verbatim}
export ANTHROPIC_API_KEY=<key>
cd clio-search
python3 eval/eval_final/code/laptop/\
  L1_three_path.py
\end{verbatim}

\artcomp

$T_1$: run \texttt{L1\_three\_path.py} --- for each of 10 queries, issue it through all three paths (Claude Native Agent, NDP-MCP, CLIO); record per-path total tokens (main + sub-agent), wall time, and answer correctness;
$T_2$: run \texttt{eval/eval\_final/code/generate\_L1\_three\_path\_plot.py} to emit Figure~2.

Dependencies: $T_1 \rightarrow T_2$.

Parameters: \texttt{claude-sonnet-4-20250514}; single repetition per (query, path) pair; CLIO indexes 341 NDP datasets locally and returns $\sim$2K tokens of ranked citations per query; the agentic loop terminates on \texttt{done} or after an invocation cap.

\artout

\texttt{eval/eval\_final/outputs/L1\_three\_path.json} containing per-query token and time breakdowns per path, and the aggregated summary (52.08\% token reduction, 91.27\% wall-time reduction). Figure~2 is emitted as \texttt{plots/fig\_L1\_three\_path.png} and \texttt{.pdf}.


%%%%%%%%%%%%%%%%%%%%%%%%%%%%%%%%%%%%%%%%%%%%%%%%%%%%%%%%
%% A3 -- IOWarp + distributed scaling (DeltaAI GH200)
%%%%%%%%%%%%%%%%%%%%%%%%%%%%%%%%%%%%%%%%%%%%%%%%%%%%%%%%

\newartifact

\artrel

Artifact $A_3$ contains the IOWarp integration experiments and the distributed scaling experiments on NCSA DeltaAI GH200 nodes. It supports $C_2$ (corpus-adaptive orchestration validated on IOWarp blobs) and $C_3$ (federated search and distributed scaling across HPC backends).

\artexp

Reproduces Table~IX (IOWarp profiling time scales sub-linearly from 17~ms at 1K blobs to 90~ms at 500K blobs for the rich-metadata variant, while the inspection rate falls from 5.0\% to 0.01\%), Figure~4 (IOWarp scaling dual-axis plot), Table~X (weak-scaling efficiency of 0.94 at 4~workers with 2.5M arXiv abstracts), and Figure~5 (distributed cross-unit probes across unit variants Pa/kPa/bar/atm and \textdegree C/\textdegree F/K return identical hit counts at sub-50~ms latency).

\arttime

Artifact Setup: 60 minutes (DeltaAI node allocation, Docker image build for IOWarp, Python environment) plus a one-time $\sim$2~hour arXiv bulk download and shard build via \texttt{slurm\_setup.sh}.
Artifact Execution: 20 hours for the 500K IOWarp rich-metadata scale; 2 hours for distributed scaling on 2.5M arXiv; 1 hour for cross-unit correctness probes.
Artifact Analysis: 10 minutes (plot generation).
Total: approximately 25 hours (dominated by the 500K IOWarp run).

\artin

\artinpart{Hardware}

NVIDIA GH200 Grace Hopper nodes (72 ARM Neoverse V2 cores, 480~GB unified CPU+GPU memory, NVLink-C2C, Slingshot-11 interconnect). The distributed scaling experiment requires 1--4 GH200 nodes with shared Lustre storage. IOWarp experiments run on a single GH200 node inside Docker.

\artinpart{Software}

\begin{itemize}
\item Linux kernel 6.x with Slurm workload manager
\item Python 3.11 with \texttt{uv}
\item Docker with the \texttt{iowarp/deps-cpu} image \url{https://hub.docker.com/u/iowarp}
\item IOWarp Chimaera runtime v1.0.3 \url{https://github.com/iowarp}
\item DuckDB 1.1.0
\item 2.5M arXiv abstract corpus (obtained via the arXiv bulk data API)
\end{itemize}

\artinpart{Datasets / Inputs}

Synthetic IOWarp blobs are generated at experiment time under four tags (temperature in \textdegree C, pressure in kPa, wind in m/s, and mixed). The arXiv abstract corpus is obtained via the public arXiv bulk data API and split into four DuckDB shards of 625K abstracts each.

\artinpart{Installation and Deployment}

One-time prep (required before any distributed experiment): \texttt{sbatch eval/eval\_final/code/delta/slurm\_setup.sh}. This creates the venv, downloads the 2.5M arXiv corpus via \texttt{download\_arxiv.py}, builds four DuckDB shards via \texttt{index\_shard.py}, and stages weak-scaling phase directories.

Then submit the provided Slurm scripts from \texttt{eval/eval\_final/code/delta/}:
\begin{itemize}
\item \texttt{slurm\_iowarp\_l2.sh} --- IOWarp 1K to 100K (rich)
\item \texttt{slurm\_iowarp\_l2\_large.sh} --- IOWarp 500K (rich)
\item \texttt{slurm\_iowarp\_l2c.sh} --- sparse-metadata 1K--100K
\item \texttt{slurm\_iowarp\_l2c\_large.sh} --- sparse-metadata 500K (optional)
\item \texttt{slurm\_strong\_scaling.sh}, \texttt{slurm\_weak\_scaling.sh} --- distributed scaling on arXiv (Table~X)
\item \texttt{slurm\_d4\_d5.sh} --- distributed cross-unit correctness (Figure~5) and NumConQ numeric-range queries
\end{itemize}

\artcomp

Workflow for IOWarp scaling (Table~IX, Figure~4):
$T_1$: generate synthetic blobs at each scale (1K, 5K, 10K, 50K, 100K, 500K);
$T_2$: index blobs into DuckDB via the IOWarp connector;
$T_3$: profile the corpus and issue four cross-unit queries (using mismatched units: \textdegree F, psi, km/h);
$T_4$: record profiling time, inspection rate, and match counts.

Workflow for distributed scaling (Table~X, Figure~5):
$T_5$: one-time setup via \texttt{slurm\_setup.sh} (downloads arXiv and builds four DuckDB shards through \texttt{download\_arxiv.py} + \texttt{index\_shard.py});
$T_6$: launch coordinator and worker processes over TCP/JSON (\texttt{distributed\_clio.py}) via \texttt{slurm\_strong\_scaling.sh} / \texttt{slurm\_weak\_scaling.sh};
$T_7$: issue 80 queries per configuration (strong and weak scaling) and record p50/p95 latency and weak-scaling efficiency (Table~X);
$T_8$: run \texttt{slurm\_d4\_d5.sh} to issue cross-unit probes (pressure and temperature) across the distributed shards and record per-unit hit counts and latency (Figure~5, \texttt{D\_cross\_unit.json}).

Dependencies: $T_1 \rightarrow T_2 \rightarrow T_3 \rightarrow T_4$ for IOWarp; $T_5 \rightarrow \{T_6 \rightarrow T_7,\ T_8\}$ for distributed experiments ($T_6$--$T_7$ and $T_8$ are independent given $T_5$).

Parameters: IOWarp rich variant uses 2.5 measurements per blob, sparse variant uses 0.75; top-$K{=}50$ for CLIO queries; distributed weak scaling uses 625K abstracts per worker; 80 queries per configuration with no repetitions (latency percentiles computed over 80 samples).

\artout

Raw results are written as JSON to \texttt{eval/eval\_final/outputs/} (files prefixed \texttt{L2\_}, \texttt{L2C\_}, and \texttt{D\_}). The analysis pipeline reads these JSON files and produces Table~IX, Table~X, Figure~4 (IOWarp scaling dual-axis plot), and Figure~5 (distributed cross-unit correctness bars). Plots are generated via \texttt{eval/eval\_final/code/generate\_plots.py}.

\end{document}
